\noindent\textbf{Elección del procedimiento.} Despreciando masa de combustible y presión residual, $F=\dot m(v_e-v_0)$; por eso menor incremento de velocidad requiere mayor caudal. La potencia cinética añadida es $\dot m(v_e^2-v_0^2)/2$, y al dividir $Fv_0$ se obtiene $\eta_p=2v_0/(v_e+v_0)$.

\begin{enumerate}

\item Despejar el gasto del turborreactor.

\[\fitmath{\mathrm{mdj}=\frac{F}{\mathrm{vej}-v_{0}}}\]

\[\fitmath{\mathrm{mdj}=\frac{(20000)}{(600)-(230)}\approx 54{,}05405\,\mathrm{kg/s}}\]

\item Evaluar su rendimiento cinético ideal.

\[\fitmath{\mathrm{etaj}=\frac{2\cdot v_{0}}{\mathrm{vej}+v_{0}}}\]

\[\fitmath{\mathrm{etaj}=\frac{2\cdot (230)}{(600)+(230)}\approx 0{,}5542169\,}\]

\item El turbofán necesita mayor gasto para el mismo empuje.

\[\fitmath{\mathrm{mdfan}=\frac{F}{\mathrm{ve}-v_{0}}}\]

\[\fitmath{\mathrm{mdfan}=\frac{(20000)}{(330)-(230)}\approx 200\,\mathrm{kg/s}}\]

\item Comparar a la misma velocidad de vuelo.

\[\fitmath{\mathrm{etafan}=\frac{2\cdot v_{0}}{\mathrm{ve}+v_{0}}}\]

\[\fitmath{\mathrm{etafan}=\frac{2\cdot (230)}{(330)+(230)}\approx 0{,}8214286\,}\]

\end{enumerate}
